% ===========================================================================
\section{Conclusions}
% ===========================================================================

Superconducting cavities for axion haloscopes must hold a high quality factor in a multi-tesla field, and copper is the substrate that makes such cavities buildable. This work took the Cu-Sn bronze route from coupons into real resonators and measured what a cavity does with it.

A hexagonal prism replaced the conventional cylinder for the prototype cavities. Simulation shows that two resistive seams cost a cylinder \qty{27}{\percent} of its quality factor and a hexagon \qty{1}{\percent}, because a hexagon concentrates surface current at the face centres and starves the corners where the seams sit. An eight-piece hexagonal prototype with one Nb$_3$Sn wall confirmed that the coupon recipe transfers to cavity scale: cavity and coupon both show $T_c = \qty{15}{\kelvin}$, and the coated wall gave $R_s = \qty{8}{\milli\ohm}$ against \qty{11}{\milli\ohm} for copper.

A two-piece hexagonal cavity, fully coated by the Nb-first Cu-Sn route, reached $Q_L = 77{,}000$ at \qty{50}{\milli\kelvin} in a dilution refrigerator, flat from \qty{50}{} to \qty{300}{\milli\kelvin}. Correcting for endcaps that lost their film during the APS etch and for leakage at the single seam gives $Q_0 \approx 1.4\times10^{5}$, a factor of \num{2.5} above the bare copper reference. This is the first Nb$_3$Sn-coated RF cavity built on a copper substrate by the Cu-Sn ternary bronze route.

Coating that same body with only its Ta/Nb underlayer --- \qty{500}{\nano\meter} Ta and \qty{750}{\nano\meter} Nb, both DC sputtered at \qty{200}{\degreeCelsius} --- reached $Q_0 = 1.4\times10^{6}$ at \qty{2}{\kelvin}, twenty-five times what the bare body gave and ten times the corrected Nb$_3$Sn value. Above the transition the same cavity falls to $Q = 13{,}000$ at \qty{9}{\kelvin}, below bare copper, because normal-state Nb is the more resistive metal and the film is thicker than its skin depth. Niobium cannot serve a \qty{9}{\tesla} haloscope, so this cavity is a control rather than a candidate. Its value is what it eliminates: the copper body, the seam, the geometry, and the measurement chain all support $10^{6}$. The order of magnitude lost between Ta/Nb and Nb$_3$Sn therefore belongs to the Cu-Sn evaporation, the \qty{715}{\degreeCelsius} reaction, and the etch, not to Nb$_3$Sn itself.

\TODO{One paragraph on the electropolished cylindrical copper baseline, once measured}

\TODO{One paragraph on the cylindrical hot-bronze Nb$_3$Sn cavity, once measured, including whether $Q(B)$ improved over the eight-piece result}

The near-term work is therefore process, not materials. Terminating the coating in Nb$_3$Sn by the hot-bronze route removes the etch that stripped the endcaps. Rotating or tilting the cavity during deposition puts film on surfaces perpendicular to the source. Electropolishing a cylindrical body raises the copper baseline that every coating is measured against. Beyond zero field, the CTE mismatch with copper remains the governing limit: it drives the film into compressive strain, lowers the vortex entry field to about \qty{11.4}{\milli\tesla}, and produces the linear $Q$ degradation and \qty{8}{\tesla} breakdown observed here. Buffer layers and multilayer architectures that decouple Nb$_3$Sn from copper stress are the route to holding $Q$ at the fields where axion searches actually run.
